\clearpage
\thispagestyle{empty}
\SourceLeaf{381}\pdfbookmark[0]{Perturbations of Ceres}{ceres}
\vspace*{1.6in}
\begin{center}{\Large PERTURBATIONS OF CERES.\par}\end{center}
\clearpage
\fancyhead[L]{\small PERTURBATIONS OF CERES}
\SourceLeaf{383}\pdfbookmark[1]{Correspondence and Posthumous Papers}{ceres-briefwechsel}
\begin{center}
{\large CORRESPONDENCE AND POSTHUMOUS PAPERS.\par}
\bigskip
[I. FIRST METHOD.]\par
[OCTOBER—DECEMBER 1802.]\par
\bigskip
\pdfbookmark[2]{1. Letters, 12 and 26 October 1802; 28 December 1815}{ceres-001}[1.]\par
\medskip
\textsc{Gauss} to \textsc{Olbers}, Braunschweig, 12 October 1802.
\end{center}
\noindent\hspace*{\parindent}.\quad.\quad.\quad.\quad.\quad I have begun working to improve the orbit of {\AstroSymbols ⚳}. First I calculated the perturbations from set VII of the elements, taking into account only the first power of the eccentricity for now, since the arc is still too small to repay the considerable labour of calculating the higher powers. I find some considerable differences from \textsc{Oriani}'s results. I have already obtained a new orbit with these perturbations and in turn redetermined the perturbations from this new orbit. In all I have 54 equations. With these perturbations, which are not \textls[180]{very} different from the first, I shall calculate a new orbit tomorrow or the day after, and then send you the result as a whole. The perturbations found in longitude are:
\begin{center}
\begingroup\footnotesize\setlength{\tabcolsep}{4pt}\renewcommand{\arraystretch}{1.25}
\begin{tabular}{l|c|c|c|c|c}
 & 1801 Jan. 1 & Feb. 11 & Dec. 7 & 1802 Mar. 17 & May 24\\
\hline
by the 1st form. & \(-428''{,}72\) & \(-358''{,}23\) & \(+646''{,}91\) & \(+954''{,}39\) & \(+1083''{,}37\)\\
»\quad »\quad 2nd\quad » & \(-415{,}82\) & \(-346{,}12\) & \(+641{,}41\) & \(+945{,}31\) & \(+1069{,}42\)\\
\end{tabular}\endgroup
\end{center}
\par
\noindent\hspace*{\parindent}In latitude the difference is quite insignificant, and I have yet to finish with the radius vector. — In the time it takes me to calculate the numerical value from the perturbation formulas, I could get more than one orbit determination \SourceLeaf{384}done. I hope, however, to be able to represent all the observations available so far almost exactly throughout, and can thus expect very good agreement even next year. — For Ceres we can probably still obtain usable convergence for the perturbations, and I have already sketched out some ideas for conveniently shortening and condensing the tables, but for {\AstroSymbols ⚴} the theory will surely still require a kind of revolution.\quad.\quad.\quad.\quad.\quad.
\par\bigskip
\begin{center}\textsc{Gauss} to \textsc{Olbers}, Braunschweig, 26 October 1802.\end{center}
\noindent\hspace*{\parindent}.\quad.\quad.\quad.\quad.\quad I have finished the calculations concerning {\AstroSymbols ⚳} mentioned in my last letter.\quad.\quad.\quad.\quad.\quad.
\par\medskip
\begin{center}\rule{.27\linewidth}{.4pt}\end{center}
\par\bigskip
\Needspace{8\baselineskip}
\begin{center}\textsc{Gauss} to \textsc{Gerling}, Göttingen, 28 December 1815.\end{center}
\noindent\hspace*{\parindent}.\quad.\quad.\quad.\quad.\quad I myself had worked out how to arrange the calculation of perturbations by the usual procedure in 1802, when I applied it to Ceres. My papers, however, are of course in such a state that I alone could use them. Nevertheless, the instructions of \textsc{Laplace} in the Mec. Cel. Vol. I, or those of \textsc{Schubert} B. 3, must agree with it in essentials.\quad.\quad.\quad.\quad.\quad.
\par\medskip
\begin{center}\rule{.27\linewidth}{.4pt}\end{center}
